\Needspace{6\baselineskip}\Article{353}
\textit{Erstes Beispiel für \(n=19\).} Da hier \(n-1=3.3.2\) gilt, ist die Bestimmung der Wurzeln \(\Omega\) auf die Lösung zweier kubischer Gleichungen und einer quadratischen Gleichung zurückzuführen. Dieses Beispiel wird umso leichter verständlich sein, als die nötigen Operationen größtenteils schon im Vorhergehenden enthalten sind. Wählt man für die primitive Wurzel \(g\) die Zahl 2, so ergeben sich die folgenden kleinsten Reste ihrer Potenzen (die Exponenten der Potenzen stehen in der ersten Zeile über den Resten): {\small\setlength{\arraycolsep}{2pt}\[\begin{array}{rrrrrrrrrrrrrrrrrr}0.&1.&2.&3.&4.&5.&6.&7.&8.&9.&10.&11.&12.&13.&14.&15.&16.&17\\1.&2.&4.&8.&16.&13.&7.&14.&9.&18.&17.&15.&11.&3.&6.&12.&5.&10\end{array}\]}

Hieraus lässt sich nach den Artikeln 344 und 345 leicht die folgende Einteilung sämtlicher Wurzeln \(\Omega\) in drei Perioden mit je sechs Gliedern und jeder einzelnen von diesen in drei Perioden mit je zwei Gliedern ableiten: \[\Omega=(18,1)\left\{\begin{array}{l}(6,1)\left\{\begin{array}{ll}(2,1)\ldots&[1],[18]\\2,8)\ldots&[8],[11]\\(2,7)\ldots&[7],[12]\end{array}\right.\\[8pt](6,2)\left\{\begin{array}{ll}(2,2)\ldots&[2],[17]\\(2,16)\ldots&[3],[16]\\(2,14)\ldots&[5],[14]\end{array}\right.\\[8pt](6,4)\left\{\begin{array}{ll}(2,4)\ldots&[4],[15]\\(2,13)\ldots&[6],[13]\\(2,9)\ldots&[9],[10]\end{array}\right.\end{array}\right.\]

Die Gleichung \((A)\), deren Wurzeln die Summen \((6,1),\allowbreak (6,2),\allowbreak (6,4)\) sind, ergibt sich als \(x^3+xx-6x-7=0\); eine ihrer Wurzeln findet man gleich \(-1{,}2218761623\). Bezeichnen wir diese mit \((6,1)\), so wird
\SourceLeaf{440}
\[\begin{aligned}(6,2)&=4-(6,1)^2=2{,}5070186441\\(6,4)&=-5-(6,1)+(6,1)^2=-2.2851424818\end{aligned}\] Damit ist \(X\) in drei Faktoren vom Grad 6 zerlegt, wenn diese Werte in Artikel 348 eingesetzt werden.

Die Gleichung \((B)\), deren Wurzeln die Summen \((2,1),\allowbreak (2.7),\allowbreak (2,8)\) sind, lautet \[x^3-(6,1)xx+((6,1)+(6,4))x-2-(6,2)=0\] oder \[x^3+1{,}2218761623xx-3{,}5070186441x-4{,}5070186441=0\] Eine ihrer Wurzeln findet man gleich \(-1{,}3545631433\), die wir mit \((2,1)\) bezeichnen werden. Nach dem Verfahren von Artikel 346 erhält man aber die folgenden Gleichungen, in denen der Kürze wegen \(q\) für \((2,1)\) geschrieben wird: \[\begin{aligned}
(2,2)&=qq-2,\quad (2,3)=q^3-3q,\\
(2,4)&=q^4-4qq+2,\quad (2,5)=q^5-5q^3+5q\\
(2,6)&=q^6-6q^4+9qq-2,\\
(2,7)&=q^7-7q^5+14q^3-7q\\
(2,8)&=q^8-8q^6+20q^4-16qq+2,\\
(2,9)&=q^9-9q^7+27q^5-30q^3+9q
\end{aligned}\]

Bequemer als nach den Regeln von Artikel 346 lassen sich diese Gleichungen im vorliegenden Fall durch die folgenden Überlegungen aufstellen. Setzt man \[[1]=\cos\frac{kP}{19}+i\sin\frac{kP}{19}\] so folgt \[\begin{aligned}
\relax[18]&=\cos\frac{18kP}{19}+i\sin\frac{18kP}{19}\\
&=\cos\frac{kP}{19}-i\sin\frac{kP}{19},
\end{aligned}\] und daher \((2,1)=2\cos\frac{kP}{19}\), sowie allgemein \[[\lambda]=\cos\frac{\lambda kP}{19}+i\sin\frac{\lambda kP}{19},\] und daher \((2,\lambda)=[\lambda]+[18\lambda]=[\lambda]+[-\lambda]=2\cos\frac{\lambda kP}{19}\) Ist also \(\frac12q=\cos\omega\), so gilt \((2,2)=2\cos2\omega,\allowbreak (2,3)=2\cos3\omega\) usw.; daraus erhält man nach den bekannten Gleichungen für die Kosinuswerte vielfacher Winkel dieselben Formeln wie oben. — Aus diesen Formeln ergeben sich nun die folgenden Zahlenwerte: \[\begin{array}{r|r}(2,2)=-0{,}1651586909&(2,6)=0{,}4909709743\\(2,3)=1{,}5782810188&(2,7)=-1{,}7589475024\\(2,4)=-1{,}9727226068&(2,8)=1{,}8916344834\\(2,5)=1{,}0938963162&(2,9)=-0{,}8033908493\end{array}\]
\SourceLeaf{441}

Die Werte von \((2,7),\allowbreak (2,8)\) können auch aus der Gleichung \((B)\) gewonnen werden, deren beide übrigen Wurzeln sie sind. Der Zweifel, welche dieser Wurzeln \((2,7)\) und welche \((2,8)\) ist, lässt sich entweder durch eine Näherungsrechnung nach den vorstehenden Formeln oder durch Sinustafeln beheben. Schon ein flüchtiger Blick in diese zeigt, dass \((2,1)=2\cos\omega\) gilt, wenn man \(\omega=\frac7{19}P\) setzt; daraus muss \[(2,7)=2\cos\frac{49}{19}P=2\cos\frac8{19}P,\] und \((2,8)=2\cos\frac{56}{19}P=2\cos\frac1{19}P\) folgen. Ebenso lassen sich die Summen \((2,2),\allowbreak (2,3),\allowbreak (2,5)\) auch durch die Gleichung \[x^3-(6,2)xx+((6,1)+(6,2))x-2-(6,4)=0\] finden, deren Wurzeln sie sind; die Ungewissheit, welche Wurzeln jeweils jenen Summen gleichzusetzen sind, wird auf genau dieselbe Weise wie zuvor beseitigt. Ebenso können auch die Summen \((2,4),\allowbreak (2,6),\allowbreak (2,9)\) durch die Gleichung \[x^3-(6,4)xx+((6,2)+(6,4))x-2-(6,1)=0\] ermittelt werden.

Schließlich sind \([1]\) und \([18]\) die Wurzeln der Gleichung \(xx-(2,1)x+1=0\), von denen die eine \(=\frac12(2,1)+i\surd(1-\frac14(2,1)^2)=\frac12(2,1)+i\surd(\frac12-\frac14(2,2))\), die andere \(=\frac12(2,1)-i\surd(\frac12-\frac14(2,2))\) ist; daraus ergeben sich die Zahlenwerte \(=-0{,}6772815716\pm0{,}7357239107i\). Die sechzehn übrigen Wurzeln können entweder durch Entwicklung der Potenzen einer beliebigen dieser beiden Wurzeln oder durch Lösung von acht anderen ähnlichen Gleichungen abgeleitet werden. Bei dem letzteren Verfahren muss man entweder durch Sinustafeln oder durch den im folgenden Beispiel zu erläuternden Kunstgriff entscheiden, bei welcher Wurzel der imaginäre Teil das positive und bei welcher das negative Vorzeichen erhalten soll. So wurden die folgenden Werte gefunden, wobei das obere Zeichen zur ersten und das untere zur zweiten Wurzel gehören soll: \[\begin{aligned}\relax[1]\text{ und }[18]&=-0{,}6772815716\pm0{,}7357239107i\\\relax[2]\text{ und }[17]&=-0{,}0825793455\mp0{,}9965844930i\\\relax[3]\text{ und }[16]&=0{,}7891405094\pm0{,}6142127127i\\\relax[4]\text{ und }[15]&=-0{,}9863613034\pm0{,}1645945903i\\\relax[5]\text{ und }[14]&=0{,}5469481581\mp0{,}8371664783i\\\relax[6]\text{ und }[13]&=0{,}2454854871\pm0{,}9694002659i\\\relax[7]\text{ und }[12]&=-0{,}8794737512\mp0{,}4759473930i\\\relax[8]\text{ und }[11]&=0{,}9458172417\mp0{,}3246994692i\\\relax[9]\text{ und }[10]&=-0.4016954247\pm0{,}9157733267i\end{aligned}\]

\SourceLeaf{442}

\Needspace{6\baselineskip}\Article{354}
\textit{Zweites Beispiel für \(n=17\).} Hier gilt \(n-1=2.2.2.2\), weshalb die Berechnung der Wurzeln \(\Omega\) auf vier quadratische Gleichungen zurückzuführen ist. Als primitive Wurzel wählen wir hier die Zahl 3, deren Potenzen die folgenden kleinsten Reste nach dem Modul 17 liefern: {\small\setlength{\arraycolsep}{2pt}\[\begin{array}{rrrrrrrrrrrrrrrr}0.&1.&2.&3.&4.&5.&6.&7.&8.&9.&10.&11.&12.&13.&14.&15\\1.&3.&9.&10.&13.&5.&15.&11.&16.&14.&8.&7.&4.&12.&2.&6\end{array}\]}

Hieraus ergeben sich die folgenden Einteilungen der Menge \(\Omega\) in zwei Perioden mit je acht, vier mit je vier und acht mit je zwei Gliedern: {\small\setlength{\arraycolsep}{2pt}\[\Omega=(16,1)\left\{\begin{array}{l}(8,1)\left\{\begin{array}{l}(4,1)\left\{\begin{array}{ll}(2,1)\ldots&[1],[16]\\(2,13)\ldots&[4],[13]\end{array}\right.\\[4pt](4,9)\left\{\begin{array}{ll}(2,9)\ldots&[8],[9]\\(2,15)\ldots&[2],[15]\end{array}\right.\end{array}\right.\\[8pt](8,3)\left\{\begin{array}{l}(4,3)\left\{\begin{array}{ll}(2,3)\ldots&[3],[14]\\(2,5)\ldots&[5],[12]\end{array}\right.\\[4pt](4.10)\left\{\begin{array}{ll}(2,10)\ldots&[7],[10]\\(2,11)\ldots&[6],[11]\end{array}\right.\end{array}\right.\end{array}\right.\]}

Die Gleichung \((A)\), deren Wurzeln die Summen \((8,1),\allowbreak (8,3)\) sind, ergibt sich nach den Regeln von Artikel 351 als \(xx+x-4=0\). Ihre Wurzeln berechnet man zu \(-\frac12+\frac12\surd 17=1{,}5615528128\) und \(-\frac12-\frac12\surd 17=-2{,}5615528128\). Die erste setzen wir \(=(8,1)\), weshalb die zweite notwendig \(=(8,3)\) zu setzen ist.

Ferner ergibt sich die Gleichung, deren Wurzeln die Summen \((4,1)\) und \((4,9)\) sind, als \((B):xx-(8,1)x-1=0\). Ihre Wurzeln sind \(\frac12(8,1)\pm\frac12\surd(4+(8,1)^2)=\frac12(8,1)\pm\frac12\surd(12+3(8,1)+4(8,3))\). Diejenige, bei der die Wurzelgröße das positive Vorzeichen erhält und deren Zahlenwert \(2{,}0494811777\) ist, setzen wir \(=(4,1)\); dann muss die andere, bei der die Wurzelgröße negativ genommen wird und deren Wert \(-0{,}4879283649\) ist, mit \((4,9)\) bezeichnet werden. Die übrigen Summen mit vier Gliedern, nämlich \((4,3)\) und \((4,10)\), können jedoch auf zwei Arten bestimmt werden. \textit{Erstens} durch das Verfahren von Artikel 346, das die folgenden Formeln liefert, in denen zur Abkürzung für \((4,1)\) der Buchstabe \(p\) geschrieben wird: \[\begin{aligned}(4,3)&=-\frac32+3p-\frac12p^3=0{,}3441507314\\(4,10)&=\frac32+2p-pp-\frac12p^3=-2{,}9057035442\end{aligned}\]
\SourceLeaf{443}

Dasselbe Verfahren liefert auch die Formel \((4,9)=-1-6p+pp+p^3\), aus der sich derselbe Wert ergibt, den wir zuvor angegeben haben. \textit{Zweitens} lassen sich die Summen \((4,3)\), \((4,10)\) auch durch Lösung der Gleichung bestimmen, deren Wurzeln sie sind. Diese Gleichung lautet \(xx-(8,3)x-1=0\), weshalb ihre Wurzeln \(\frac12(8,3)\pm\frac12\surd(4+(8,3)^2)\) oder \(\frac12(8,3)+\frac12\surd(12+4(8,1)+3(8,3))\) und \(\frac12(8,3)-\frac12\surd(12+4(8,1)+3(8,3))\) sind. Der Zweifel, \textit{welche} Wurzel mit \((4,3)\) und welche mit \((4,10)\) bezeichnet werden muss, wird durch den folgenden Kunstgriff beseitigt, den wir in Artikel 352 erwähnt haben. Man entwickle das Produkt aus \((4,1)-(4,9)\) und \((4,3)-(4,10)\); dabei ergibt sich \(2(8,1)-2(8,3)\)\HistoricalNote{*)}{Das eigentliche Wesen dieses Kunstgriffs besteht darin, dass man a priori vorhersehen konnte, dass das entwickelte Produkt keine Summen mit vier Gliedern enthalten würde, sondern sich allein durch Summen mit acht Gliedern darstellen ließe. Den Grund dafür, der hier der Kürze wegen übergangen werden muss, werden mit dem Gegenstand Vertraute sehr leicht erkennen.} Der Wert dieses Ausdrucks ist nun offenbar positiv, nämlich \(=+2\surd 17\), und auch der erste Faktor des Produkts, \((4,1)-(4,9)\), ist positiv, nämlich \(=+\surd(12+3(8,1)+4(8,3))\). Daher muss notwendig auch der andere Faktor \((4,3)-(4,10)\) positiv sein, und folglich ist \((4,3)\) der \textit{ersten} Wurzel, bei der das positive Vorzeichen vor der Wurzelgröße steht, und \((4,10)\) der zweiten gleichzusetzen. Hieraus ergeben sich dieselben Zahlenwerte wie oben.

Nachdem sämtliche Summen mit vier Gliedern gefunden sind, gehen wir zu den Summen mit zwei Gliedern über. Die Gleichung \((C)\), deren Wurzeln \((2,1),\allowbreak (2,13)\) sind, die in \((4,1)\) enthalten sind, ergibt sich als \(xx-(4,1)x+(4,3)=0\). Ihre Wurzeln sind \(\frac12(4,1)\pm\frac12\surd(-4(4,3)+(4,1)^2)\) oder \(\frac12(4,1)\pm\frac12\surd(4+(4,9)-2(4,3))\). Diejenige, bei der die Wurzelgröße positiv genommen wird und deren Wert sich als \(=1{,}8649444588\) ergibt, setzen wir \(=(2,1)\); dann wird \((2,13)\) der anderen gleich, deren Wert \(=0{,}1845367189\) ist. Will man die übrigen Summen mit zwei Gliedern nach dem Verfahren von Artikel 346 bestimmen, so können für \((2,2),\allowbreak (2,3),\allowbreak (2,4),\allowbreak (2,5),\allowbreak (2,6),\allowbreak (2,7),\allowbreak (2,8)\) dieselben Formeln verwendet werden, die wir im vorigen Beispiel für die ebenso bezeichneten Größen angegeben haben, nämlich \((2,2)\), (oder \((2,15)\)), \(=(2,1)^2-2\) usw. Will man sie dagegen lieber paarweise durch Lösung einer quadratischen Gleichung berechnen, so findet man für \((2,9),\allowbreak (2,15)\) die Gleichung \(xx-(4,9)x+(4,10)=0\), deren Wurzeln \(\frac12(4,9)\pm\frac12\surd(4+(4,1)-2(4,10))\) sind. Wie hier über das doppelte Vorzeichen zu entscheiden ist, lässt sich auf dieselbe Weise wie oben bestimmen. Entwickelt man nämlich das Produkt aus \((2,1)-(2,13)\) und \((2,9)-(2,15)\), so erhält man \(-(4,1)+(4,9)-(4,3)+(4,10)\). Da dies offenbar negativ, der Faktor \((2,1)-(2,13)\) aber positiv ist, muss notwendig \((2,9)-(2,15)\)
\SourceLeaf{444}
negativ sein; daher ist in dem zuvor angegebenen Ausdruck für \((2,15)\) das obere, positive und für \((2,9)\) das untere, negative Vorzeichen zu wählen. Daraus berechnet man \((2,9)=-1{,}9659461994\), \((2,15)=1{,}4780178344\). — Da ebenso aus der Entwicklung des Produkts aus \((2,1)-(2,13)\) und \((2,3)-(2,5)\) der Ausdruck \((4,9)-(4,10)\) und damit eine positive Größe hervorgeht, schließen wir, dass der Faktor \((2,3)-(2,5)\) positiv ist. Eine ähnliche Rechnung wie zuvor liefert dann \[\begin{aligned}(2,3)&=\tfrac12(4,3)+\tfrac12\surd(4+(4,10)-2(4,9))\\&=0{,}8914767116\\(2,5)&=\tfrac12(4,3)-\tfrac12\surd(4+(4,10)-2(4,9))\\&=-0{,}5473259801\end{aligned}\]

Schließlich erhält man durch ganz entsprechende Operationen \[\begin{aligned}(2,10)&=\tfrac12(4,10)-\tfrac12\surd(4+(4,3)-2(4,1))\\&=-1{,}7004342715\\(2,11)&=\tfrac12(4,10)+\tfrac12\surd(4+(4,3)-2(4,1))\\&=-1{,}2052692728\end{aligned}\]

Es bleibt, zu den Wurzeln \(\Omega\) selbst herabzusteigen. Die Gleichung \((D)\), deren Wurzeln \([1]\) und \([16]\) sind, lautet \(xx-(2,1)x+1=0\), woraus sich die Wurzeln \(\frac12(2,1)\pm\frac12\surd((2,1)^2-4)\) oder vielmehr \(\frac12(2,1)\pm i\surd(4-(2,1)^2)\) oder \(\frac12(2,1)\pm\frac12i\surd(2-(2,15))\) ergeben. Für \([1]\) wählen wir das obere, für \([16]\) das untere Vorzeichen. Die vierzehn übrigen Wurzeln erhält man entweder durch die Potenzen von \([1]\) oder durch Lösung von sieben quadratischen Gleichungen, von denen jede zwei Wurzeln liefert; dabei lässt sich die Ungewissheit über die Vorzeichen der Wurzelgrößen durch denselben Kunstgriff wie zuvor beseitigen. So sind \([4]\) und \([13]\) die Wurzeln der Gleichung \(xx-(2,13)x+1=0\) und daher \(\frac12(2,13)\pm\frac12i\surd(2-(2,9))\). Aus der Entwicklung des Produkts aus \([1]-[16]\) und \([4]-[13]\) ergibt sich aber \((2,5)-(2,3)\) und damit eine negative reelle Größe. Da nun \([1]-[16]\) gleich \(+i\surd(2-(2,15))\) ist, also das Produkt der imaginären Größe \(i\) mit einer \textit{positiven} reellen Größe, muss auch \([4]-[13]\) das Produkt von \(i\) mit einer \textit{positiven} reellen Größe sein, da \(ii=-1\) gilt. Daraus folgt, dass für \([4]\) das obere und für \([13]\) das untere Vorzeichen zu nehmen ist. Ebenso findet man für die Wurzeln \([8]\) und \([9]\) den Ausdruck \(\frac12(2,9)\pm\frac12i\surd(2-(2,1))\). Da das Produkt aus \([1]-[16]\) und \([8]-[9]\) gleich \((2,9)-(2,10)\) und daher negativ ist, muss für \([8]\) das obere und für \([9]\) das untere Vorzeichen gewählt werden. Berechnen wir die übrigen Wurzeln auf dieselbe Weise, so erhalten wir die folgenden Zahlenwerte, wobei die oberen Vorzeichen zu den ersten und die unteren zu den zweiten Wurzeln gehören sollen:
\SourceLeaf{445}
\[\begin{aligned}\relax[1],\ [16]&\quad\ldots\quad 0{,}9324722294\pm 0{,}3612416662i\\\relax [2],\ [15]&\quad\ldots\quad 0{,}7390089172\pm 0{,}6736956436i\\\relax [3],\ [14]&\quad\ldots\quad 0{,}4457383558\pm 0{,}8951632914i\\\relax [4],\ [13]&\quad\ldots\quad 0{,}0922683595\pm 0{,}9957341763i\\\relax [5],\ [12]&\quad\ldots\quad -0{,}2736629901\pm 0{,}9618256432i\\\relax [6],\ [11]&\quad\ldots\quad -0{,}6026346364\pm 0{,}7980172273i\\\relax [7],\ [10]&\quad\ldots\quad -0{,}8502171357\pm 0{,}5264321629i\\\relax [8],\ [9]&\quad\ldots\quad -0{,}9829730997\pm 0{,}1837495178i\end{aligned}\] \[\begin{matrix}*\qquad *\\{}*\end{matrix}\]

Das Vorstehende könnte zwar zur Lösung der Gleichung \(x^n-1=0\) und damit auch zur Bestimmung der trigonometrischen Funktionen ausreichen, die zu Bogen gehören, welche mit dem Umfang kommensurabel sind. Wegen der Bedeutung des Gegenstands können wir diese Untersuchung jedoch nicht abschließen, ohne zuvor aus der Fülle sowohl der ihn erläuternden Beobachtungen als auch der mit ihm verwandten oder von ihm abhängigen Sätze einiges hier anzufügen. Dabei werden wir vor allem solche auswählen, die sich ohne einen umfangreichen Apparat anderer Untersuchungen erledigen lassen; wir wollen sie nur als \textit{Proben} dieser sehr umfassenden, künftig ausführlich zu behandelnden Theorie betrachtet wissen.

